% Optional one-minute bridge; enable the indicated line in main.tex.
\begin{frame}[t]{GLMs: a brief connection}
\lead{A generalized linear model links a conditional mean to a linear predictor.}
\begin{definitionbox}
\[\mu(x)=\E[Y\mid x],\qquad g(\mu(x))=\beta_0+x^\top\beta.\]
A conditional exponential-family distribution specifies the likelihood.
\end{definitionbox}
\vspace{.1cm}
\centering\small\renewcommand{\arraystretch}{1.2}
\begin{tabular}{@{}lll@{}}\toprule
Response model & Mean & Link $g(\mu)$\\\midrule
Gaussian & real-valued & $\mu$\\
Bernoulli & probability & $\log\!\frac{\mu}{1-\mu}$\\
Poisson & positive count mean & $\log\mu$\\\bottomrule
\end{tabular}
\vspace{.15cm}
\begin{takeaway}For ensembles, our next step is to learn the predictor as a sum of simple functions.\end{takeaway}
\vfill{\scriptsize\href{https://web.stanford.edu/class/stats305b/notes/GLM_I.html}{Taylor, Stanford STATS 305B: Generalized linear models.}}
\note{This is optional background only and takes about one minute. The displayed links are canonical examples, not a claim that every GLM must use its canonical link. A GLM specifies distribution, predictor and link. The link acts on the conditional mean, not directly on the observed response. No exponential-family derivation is needed for the subsequent ensemble material. In this lecture the main path runs directly from the regression recap to ensembles, without this slide.}
\end{frame}
